%!TEX root = ../Thesis.tex
\chapter{Notation}

\begin{table}[H]
\centering
\begin{tabular}{r p{10cm}}
	symbol & meaning \\ \hline
	$a$ & The activation input, calculated as a weighted sum over the input. \\
	$b$ & The activation output, $b = \theta(a)$. \\
	$H$ & The amount of units in the layer.\\
	$h$ & The index of a units in the layer. \\
	$K$ & The amount of classes calculated in the softmax, $K = H_{L+1}$. \\ 
	$k$ & The class index used in the softmax, $k = h_{L + 1}$.  \\
	$\mathcal{L}$ & The loss function (should always be minimized). \\
	$L$ & The amount of hidden layers. Thus including the softmax output layer there are $L+1$ layers. \\
	$\ell$ & The layer index. \\
	$m$ & The hyperparameter momentum used in gradient descent. \\
	$s$ & The cell state used in a LSTM unit. \\
	$t$ & A target value. If used in a superscript it is the sequence index. \\
	$w$ & A weight used in a neural neural network. \\
	$x$ & The input to the neural network, $x = b_{h_0}$. \\
	$y$ & The softmax output. \\
	$\gamma$ & Decay rate used in RMSprop gradient descent. \\
	$\delta$ & Bookkeeping value used in backward propagation. \\
	$\eta$ & Learning rate used in gradient descent. \\
	$\rho$ & Indicates the input gate in a LSTM unit. \\ 
	$\phi$ & Indicates the forget gate in a LSTM unit. \\ 
	$\omega$ & Indicates the output gate in a LSTM unit. \\ 
	$\theta$ & The non-linear activation function.
\end{tabular}
\caption*{\textbf{Table C:} Meaning of commonly used symbols.}
\end{table}

\vspace{-0.1cm}
Note that subscript and superscript may be added to indicate the layer and time step.
\begin{figure}[H]
	\vspace{-0.2cm}
	\centering
	\includegraphics[scale=0.7]{appendices/notation}
\end{figure}
